\documentclass[11pt,a4paper]{article}\usepackage[]{graphicx}\usepackage[]{xcolor}
% maxwidth is the original width if it is less than linewidth
% otherwise use linewidth (to make sure the graphics do not exceed the margin)
\makeatletter
\def\maxwidth{ %
  \ifdim\Gin@nat@width>\linewidth
    \linewidth
  \else
    \Gin@nat@width
  \fi
}
\makeatother

\definecolor{fgcolor}{rgb}{0.345, 0.345, 0.345}
\newcommand{\hlnum}[1]{\textcolor[rgb]{0.686,0.059,0.569}{#1}}%
\newcommand{\hlstr}[1]{\textcolor[rgb]{0.192,0.494,0.8}{#1}}%
\newcommand{\hlcom}[1]{\textcolor[rgb]{0.678,0.584,0.686}{\textit{#1}}}%
\newcommand{\hlopt}[1]{\textcolor[rgb]{0,0,0}{#1}}%
\newcommand{\hlstd}[1]{\textcolor[rgb]{0.345,0.345,0.345}{#1}}%
\newcommand{\hlkwa}[1]{\textcolor[rgb]{0.161,0.373,0.58}{\textbf{#1}}}%
\newcommand{\hlkwb}[1]{\textcolor[rgb]{0.69,0.353,0.396}{#1}}%
\newcommand{\hlkwc}[1]{\textcolor[rgb]{0.333,0.667,0.333}{#1}}%
\newcommand{\hlkwd}[1]{\textcolor[rgb]{0.737,0.353,0.396}{\textbf{#1}}}%
\let\hlipl\hlkwb

\usepackage{framed}
\makeatletter
\newenvironment{kframe}{%
 \def\at@end@of@kframe{}%
 \ifinner\ifhmode%
  \def\at@end@of@kframe{\end{minipage}}%
  \begin{minipage}{\columnwidth}%
 \fi\fi%
 \def\FrameCommand##1{\hskip\@totalleftmargin \hskip-\fboxsep
 \colorbox{shadecolor}{##1}\hskip-\fboxsep
     % There is no \\@totalrightmargin, so:
     \hskip-\linewidth \hskip-\@totalleftmargin \hskip\columnwidth}%
 \MakeFramed {\advance\hsize-\width
   \@totalleftmargin\z@ \linewidth\hsize
   \@setminipage}}%
 {\par\unskip\endMakeFramed%
 \at@end@of@kframe}
\makeatother

\definecolor{shadecolor}{rgb}{.97, .97, .97}
\definecolor{messagecolor}{rgb}{0, 0, 0}
\definecolor{warningcolor}{rgb}{1, 0, 1}
\definecolor{errorcolor}{rgb}{1, 0, 0}
\newenvironment{knitrout}{}{} % an empty environment to be redefined in TeX

\usepackage{alltt}
\usepackage[utf8]{inputenc}
\usepackage[T1]{fontenc}
\usepackage[margin=1in]{geometry}
\usepackage{graphicx}
\usepackage{fancyhdr}
\usepackage{booktabs}
\usepackage{array}
\usepackage{longtable}
\usepackage{enumitem}
\usepackage{xcolor}
\usepackage{titlesec}
\usepackage{hyperref}
\usepackage{chemformula}

% Header and footer
\pagestyle{fancy}
\fancyhf{}
\fancyhead[L]{SOP-TOC-SSM-001}
\fancyhead[R]{Revision 1.0}
\fancyfoot[C]{\thepage}
\renewcommand{\headrulewidth}{0.4pt}
\renewcommand{\footrulewidth}{0.4pt}

% Section formatting
\titleformat{\section}
  {\normalfont\Large\bfseries}{\thesection}{1em}{}
\titleformat{\subsection}
  {\normalfont\large\bfseries}{\thesubsection}{1em}{}
\titleformat{\subsubsection}
  {\normalfont\normalsize\bfseries}{\thesubsubsection}{1em}{}

% Hyperref setup
\hypersetup{
    colorlinks=true,
    linkcolor=blue,
    filecolor=magenta,      
    urlcolor=cyan,
    pdftitle={SOP: Shimadzu TOC Analyzer with SSM},
    pdfpagemode=FullScreen,
}
\IfFileExists{upquote.sty}{\usepackage{upquote}}{}
\begin{document}

% Title page
\begin{titlepage}
    \centering
    \vspace*{2cm}
    
    {\Huge\bfseries Standard Operating Procedure\par}
    \vspace{1cm}
    {\Large Shimadzu TOC Analyzer with SSM for\\Soil Samples with Varying Carbonate Concentrations\par}
    \vspace{2cm}
    
    \begin{tabular}{ll}
        \toprule
        \textbf{Document Number:} & SOP-TOC-SSM-001 \\
        \textbf{Revision:} & 1.0 \\
        \textbf{Effective Date:} & October 2025 \\
        \textbf{Review Date:} & October 2026 \\
        \bottomrule
    \end{tabular}
    
    \vfill
    
    {\large \textbf{Document Control:} This is a controlled document. Printed copies are uncontrolled and may not be current.\par}
    
\end{titlepage}

% Table of contents
\tableofcontents
\newpage

\section{Purpose and Scope}

This SOP describes the procedure for determining Total Carbon (TC), Total Inorganic Carbon (TIC), and Total Organic Carbon (TOC) in soil samples using the Shimadzu TOC Analyzer equipped with a Solid Sample Module (SSM). This procedure is applicable to soil samples with high ($>10\%$), medium (2--10\%), and low ($<2\%$) carbonate concentrations.

\section{Principle}

The Shimadzu TOC-SSM system uses high-temperature combustion (900°C) to convert all carbon in the sample to \ch{CO2}, which is then quantified by NDIR (Non-Dispersive Infrared) detection. For carbonate-bearing soils:

\begin{itemize}[noitemsep]
    \item \textbf{TC (Total Carbon)} = measured by direct combustion
    \item \textbf{TIC (Total Inorganic Carbon)} = measured by acid treatment (200°C)
    \item \textbf{TOC (Total Organic Carbon)} = TC $-$ TIC
\end{itemize}

\section{Safety and Precautions}

\subsection{Personal Protective Equipment (PPE)}
\begin{itemize}[noitemsep]
    \item Safety glasses or goggles
    \item Laboratory coat
    \item Heat-resistant gloves when handling hot ceramic boats
    \item Nitrile gloves when handling acids
\end{itemize}

\subsection{Chemical Hazards}
\begin{itemize}[noitemsep]
    \item Phosphoric acid (25\% solution): corrosive, causes severe burns
    \item High-purity oxygen: oxidizer, keep away from combustibles
    \item Soil samples: may contain pathogens or heavy metals
\end{itemize}

\subsection{Equipment Hazards}
\begin{itemize}[noitemsep]
    \item Combustion furnace reaches 900°C
    \item Acid furnace reaches 200°C
    \item Allow adequate cooling time before handling boats
\end{itemize}

\section{Equipment and Materials}

\subsection{Instrumentation}
\begin{itemize}[noitemsep]
    \item Shimadzu TOC Analyzer with Solid Sample Module (SSM)
    \item High-purity oxygen supply (99.999\%)
    \item Carrier gas supply (oxygen or nitrogen, 99.999\%)
    \item Computer with TOC Control Software
    \item Analytical balance (0.0001 g precision)
\end{itemize}

\subsection{Consumables}
\begin{itemize}[noitemsep]
    \item Ceramic sample boats (pre-combusted)
    \item 25\% phosphoric acid solution (prepared fresh weekly)
    \item Glucose or sucrose standard (certified reference material)
    \item Calcium carbonate standard (reagent grade, dried at 105°C)
\end{itemize}

\subsection{Sample Preparation Equipment}
\begin{itemize}[noitemsep]
    \item Mortar and pestle or ball mill
    \item 2 mm and 0.25 mm sieves
    \item Drying oven (105°C)
    \item Desiccator with desiccant
\end{itemize}

\section{Reagents and Standards}

\subsection{Phosphoric Acid Solution (25\% v/v)}
\begin{enumerate}[noitemsep]
    \item Add 250 mL concentrated phosphoric acid (85\%) to 600 mL deionized water
    \item Mix thoroughly and dilute to 1000 mL
    \item Store in glass bottle, prepare fresh weekly
\end{enumerate}

\subsection{TOC Calibration Standards}
Prepare from glucose or sucrose:
\begin{itemize}[noitemsep]
    \item Stock solution: 1000 mg C/L
    \item Working standards: 10, 25, 50, 100, 250, 500 mg C/L
    \item Prepare fresh daily
\end{itemize}

\subsection{TIC Calibration Standards}
Prepare from dried calcium carbonate (\ch{CaCO3}):
\begin{itemize}[noitemsep]
    \item Calculate carbon content: 12.0\% C by weight
    \item Prepare solid standards: 1, 5, 10, 25, 50 mg \ch{CaCO3}
    \item Store in desiccator
\end{itemize}

\section{Sample Preparation}

\subsection{Initial Processing}
\begin{enumerate}[noitemsep]
    \item Air-dry soil samples at room temperature or 40°C maximum
    \item Remove large debris, roots, and stones
    \item Grind samples using mortar and pestle or ball mill
    \item Pass through 2 mm sieve for initial homogenization
\end{enumerate}

\subsection{Fine Grinding}
\begin{enumerate}[noitemsep]
    \item Take representative subsample (approximately 50 g)
    \item Further grind to pass through 0.25 mm sieve (60 mesh)
    \item Mix thoroughly to ensure homogeneity
\end{enumerate}

\subsection{Moisture Determination}
\begin{enumerate}[noitemsep]
    \item Weigh 2--5 g of ground sample in pre-weighed crucible
    \item Dry at 105°C for 24 hours
    \item Cool in desiccator and reweigh
    \item Calculate moisture content for dry weight correction
\end{enumerate}

\subsection{Sample Weighing Guidelines}

\begin{longtable}{p{4cm}p{3cm}p{3cm}p{3cm}}
\toprule
\textbf{Soil Type} & \textbf{TC Analysis} & \textbf{TIC Analysis} & \textbf{Expected TIC} \\
\midrule
\endfirsthead
\multicolumn{4}{c}%
{\tablename\ \thetable\ -- \textit{Continued from previous page}} \\
\toprule
\textbf{Soil Type} & \textbf{TC Analysis} & \textbf{TIC Analysis} & \textbf{Expected TIC} \\
\midrule
\endhead
\midrule
\multicolumn{4}{r}{\textit{Continued on next page}} \\
\endfoot
\bottomrule
\endlastfoot
High Carbonate ($>10\%$) & 10--20 mg & 5--10 mg & $>1.2\%$ C \\
Medium Carbonate (2--10\%) & 20--40 mg & 10--30 mg & 0.24--1.2\% C \\
Low Carbonate ($<2\%$) & 40--80 mg & 30--80 mg & $<0.24\%$ C \\
\end{longtable}

\section{Instrument Setup and Calibration}

\subsection{Daily Startup}
\begin{enumerate}[noitemsep]
    \item Turn on carrier gas and oxygen supply
    \item Check gas pressures (typically 0.4--0.5 MPa)
    \item Power on TOC analyzer and allow 30-minute warm-up
    \item Launch TOC Control Software
    \item Set furnace temperatures:
    \begin{itemize}
        \item TC furnace: 900°C
        \item TIC furnace: 200°C
    \end{itemize}
    \item Verify baseline stability ($<5$ ppm drift)
\end{enumerate}

\subsection{System Blank Check}
\begin{enumerate}[noitemsep]
    \item Analyze 3 blank ceramic boats
    \item Average blank should be $<0.1$ mg C
    \item If blank $>0.1$ mg C, clean boats or replace
\end{enumerate}

\subsection{TC Calibration}
\begin{enumerate}[noitemsep]
    \item Analyze glucose/sucrose standards in triplicate
    \item Use minimum 5 concentration levels
    \item Generate calibration curve ($r^2 > 0.999$ required)
    \item Verify linearity across working range
\end{enumerate}

\subsection{TIC Calibration}
\begin{enumerate}[noitemsep]
    \item Weigh \ch{CaCO3} standards directly into ceramic boats
    \item Add 150 $\mu$L of 25\% phosphoric acid
    \item Analyze in triplicate at each concentration
    \item Generate calibration curve ($r^2 > 0.998$ required)
\end{enumerate}

\subsection{Quality Control Standards}
\begin{enumerate}[noitemsep]
    \item Analyze certified reference soil (if available)
    \item Expected recovery: 90--110\%
    \item If outside range, recalibrate system
\end{enumerate}

\section{Analysis Procedure}

\subsection{Sample Loading}
\begin{enumerate}[noitemsep]
    \item Pre-clean ceramic boats by combusting at 900°C
    \item Cool boats in desiccator for minimum 30 minutes
    \item Tare balance with empty boat
    \item Add appropriate sample amount according to carbonate level
    \item Record exact sample weight (to 0.0001 g)
    \item Place boat in SSM autosampler carousel
\end{enumerate}

\subsection{TC Analysis}
\begin{enumerate}[noitemsep]
    \item Create sample sequence in software
    \item Enter sample IDs and weights
    \item Select TC analysis method:
    \begin{itemize}
        \item Combustion temperature: 900°C
        \item Combustion time: 3--5 minutes
        \item Oxygen flow: per manufacturer specifications
    \end{itemize}
    \item Start analysis sequence
    \item System automatically combusts each sample
    \item Software calculates TC based on calibration curve
\end{enumerate}

\subsection{TIC Analysis}
\begin{enumerate}[noitemsep]
    \item Prepare new sample boats with same soil samples
    \item Weigh samples (use same amounts as TC analysis)
    \item Add 150 $\mu$L of 25\% phosphoric acid to each boat
    \item Wait 30 seconds for initial reaction
    \item Load boats into TIC furnace position
    \item Select TIC analysis method:
    \begin{itemize}
        \item Temperature: 200°C
        \item Analysis time: 5--7 minutes
        \item Oxygen/carrier gas flow: per manufacturer specifications
    \end{itemize}
    \item Start analysis sequence
    \item Acid converts carbonates to \ch{CO2} at low temperature
    \item Software calculates TIC based on calibration curve
\end{enumerate}

\subsection{TOC Calculation}
\begin{equation}
\text{TOC (\%)} = \text{TC (\%)} - \text{TIC (\%)}
\end{equation}

Software typically calculates automatically.

\section{Quality Assurance/Quality Control}

\subsection{Replicate Analysis}
\begin{itemize}[noitemsep]
    \item Analyze each sample in triplicate minimum
    \item Calculate mean, standard deviation, and RSD
    \item Acceptable RSD: $<10\%$ for samples $>0.5\%$ C
    \item Acceptable RSD: $<15\%$ for samples $<0.5\%$ C
\end{itemize}

\subsection{Calibration Verification}
\begin{itemize}[noitemsep]
    \item Analyze calibration check standard every 10 samples
    \item Recovery must be 90--110\%
    \item If outside range, recalibrate and re-analyze affected samples
\end{itemize}

\subsection{Method Blanks}
\begin{itemize}[noitemsep]
    \item Analyze blank boat every 10 samples
    \item Blank should be $<10\%$ of lowest sample concentration
\end{itemize}

\subsection{Duplicate Samples}
\begin{itemize}[noitemsep]
    \item Analyze 10\% of samples as duplicates
    \item RPD (Relative Percent Difference) should be $<15\%$
\end{itemize}

\subsection{Certified Reference Materials}
\begin{itemize}[noitemsep]
    \item Analyze CRM at beginning and end of each run
    \item Document recovery and control chart results
\end{itemize}

\section{Carbonate-Specific Considerations}

\subsection{High Carbonate Soils ($>10\%$)}

\subsubsection{Challenges:}
\begin{itemize}[noitemsep]
    \item Vigorous reaction with acid
    \item Potential for sample loss or splashing
    \item May exceed detector range
\end{itemize}

\subsubsection{Solutions:}
\begin{itemize}[noitemsep]
    \item Use smaller sample sizes (5--10 mg for TIC)
    \item Add acid slowly in 50 $\mu$L increments
    \item Wait 1--2 minutes between acid additions
    \item Dilute samples if necessary
    \item Consider pre-treatment to remove excess carbonates for TOC determination
\end{itemize}

\subsection{Medium Carbonate Soils (2--10\%)}

\subsubsection{Considerations:}
\begin{itemize}[noitemsep]
    \item Standard protocol usually sufficient
    \item Monitor for complete carbonate dissolution
    \item Ensure adequate acid addition (150--200 $\mu$L)
    \item Extend TIC analysis time if needed (7--10 minutes)
\end{itemize}

\subsection{Low Carbonate Soils ($<2\%$)}

\subsubsection{Considerations:}
\begin{itemize}[noitemsep]
    \item Use larger sample sizes to improve detection
    \item TIC may be near detection limit
    \item Verify blank contamination is minimal
    \item Multiple replicates recommended ($n=5$)
    \item TOC may be determined by direct method after acid pre-treatment
\end{itemize}

\subsection{Acid Pre-treatment for Direct TOC (Optional)}

For high carbonate soils, direct TOC determination:
\begin{enumerate}[noitemsep]
    \item Weigh 0.5--1 g soil into beaker
    \item Add 10 mL of 1M HCl dropwise
    \item Allow to react until effervescence stops
    \item Repeat acid addition if necessary
    \item Rinse with deionized water (3$\times$)
    \item Dry at 60°C
    \item Grind and analyze for TOC directly
    \item \textbf{Note:} This method may result in loss of volatile organic compounds
\end{enumerate}

\section{Data Recording and Calculations}

\subsection{Raw Data}
Record in laboratory notebook:
\begin{itemize}[noitemsep]
    \item Sample ID and description
    \item Sample weight (dry weight basis)
    \item Analysis date and time
    \item Instrument ID and calibration date
    \item Peak area/height
    \item Calculated concentrations
\end{itemize}

\subsection{Calculations}

\subsubsection{Dry weight correction:}
\begin{equation}
\text{Dry weight} = \text{Wet weight} \times \left(1 - \frac{\text{Moisture \%}}{100}\right)
\end{equation}

\subsubsection{Carbon concentration (\%):}
\begin{equation}
\text{C (\%)} = \frac{\text{C measured (mg)}}{\text{Sample dry weight (mg)}} \times 100
\end{equation}

\subsubsection{TOC calculation:}
\begin{equation}
\text{TOC (\%)} = \text{TC (\%)} - \text{TIC (\%)}
\end{equation}

\subsubsection{Relative Percent Difference (RPD):}
\begin{equation}
\text{RPD} = \frac{|\text{Sample 1} - \text{Sample 2}|}{\left(\frac{\text{Sample 1} + \text{Sample 2}}{2}\right)} \times 100
\end{equation}

\subsection{Reporting}
Report results to appropriate significant figures:
\begin{itemize}[noitemsep]
    \item TC, TIC, TOC: 2 decimal places
    \item Include standard deviation or RSD
    \item Report on dry weight basis
    \item Note any dilutions or pre-treatments
\end{itemize}

\section{Maintenance and Troubleshooting}

\subsection{Daily Maintenance}
\begin{itemize}[noitemsep]
    \item Check gas pressures and flows
    \item Verify furnace temperatures
    \item Clean sample inlet if needed
    \item Check for leaks in gas lines
\end{itemize}

\subsection{Weekly Maintenance}
\begin{itemize}[noitemsep]
    \item Clean combustion tubes
    \item Replace quartz wool if discolored
    \item Check and clean NDIR detector cell
    \item Verify calibration standards
\end{itemize}

\subsection{Monthly Maintenance}
\begin{itemize}[noitemsep]
    \item Replace oxygen scrubber cartridge
    \item Check all O-rings and seals
    \item Clean or replace furnace tubes if necessary
    \item Perform full system calibration
\end{itemize}

\subsection{Common Issues and Solutions}

\begin{longtable}{p{3.5cm}p{4cm}p{6cm}}
\toprule
\textbf{Problem} & \textbf{Possible Cause} & \textbf{Solution} \\
\midrule
\endfirsthead
\multicolumn{3}{c}%
{\tablename\ \thetable\ -- \textit{Continued from previous page}} \\
\toprule
\textbf{Problem} & \textbf{Possible Cause} & \textbf{Solution} \\
\midrule
\endhead
\midrule
\multicolumn{3}{r}{\textit{Continued on next page}} \\
\endfoot
\bottomrule
\endlastfoot
High blank values & Contaminated boats & Pre-combust boats at 900°C \\
Poor peak shape & Incomplete combustion & Increase combustion time/temperature \\
Low TIC recovery & Insufficient acid & Increase acid volume or add incrementally \\
Carryover between samples & Incomplete purging & Extend purge time between samples \\
Baseline drift & Detector contamination & Clean NDIR cell \\
Inconsistent replicates & Sample heterogeneity & Improve grinding and mixing \\
\end{longtable}

\section{Waste Disposal}

\begin{itemize}[noitemsep]
    \item Acid waste: Neutralize with sodium bicarbonate, dispose per institutional guidelines
    \item Ceramic boats: Reuse after cleaning or dispose as non-hazardous solid waste
    \item Soil samples: Dispose per institutional chemical waste procedures
\end{itemize}

\section{References}

\begin{enumerate}
    \item Shimadzu TOC Analyzer SSM-5000A User Manual
    \item ASTM D4129 - Standard Test Method for Total and Organic Carbon in Water
    \item ISO 10694 - Soil Quality - Determination of Organic and Total Carbon after Dry Combustion
    \item Nelson, D.W. and Sommers, L.E. (1996) Total Carbon, Organic Carbon, and Organic Matter. Methods of Soil Analysis Part 3
\end{enumerate}

\section{Revision History}

\begin{table}[h]
\centering
\begin{tabular}{cccc}
\toprule
\textbf{Version} & \textbf{Date} & \textbf{Changes} & \textbf{Author} \\
\midrule
1.0 & October 2025 & Initial release & \\
\bottomrule
\end{tabular}
\end{table}

\section{Approval}

\begin{table}[h]
\centering
\begin{tabular}{cccc}
\toprule
\textbf{Role} & \textbf{Name} & \textbf{Signature} & \textbf{Date} \\
\midrule
Author & & & \\[0.5cm]
Reviewer & & & \\[0.5cm]
Approver & & & \\[0.5cm]
\bottomrule
\end{tabular}
\end{table}

\end{document}
